% !TeX root = ../navier-stokes-manuscript.tex
% Archived before the 2026-08-17 Section 2 rewrite.
\subsection{The finite adjacent rectangle}\label{sub:BodyFiniteAdjacentRectangles}

A finite list of mixed derivatives does not automatically have a value at the endpoint.
To trace one temporal row \(V_n\) strongly in \(H^M\), the equation has to control that row on both sides of the desired Sobolev height.
The reason is visible in one identity.

Let \(A=(1-\Delta)^{1/2}\), let \(s\in\R\), and begin with a smooth Hilbert-valued function \(v\).
For \(0\leq a<b\leq T\),
\begin{equation}\label{eq:BodyAdjacentTraceIdentity}
  \norm{v(b)}_{H^s}^2-
  \norm{v(a)}_{H^s}^2
  =2\operatorname{Re}
  \int_a^b
  \pair{A^{s-1}v_t(t)}{A^{s+1}v(t)}_{L^2}\,dt.
\end{equation}
The product on the right is integrable when
\begin{equation}\label{eq:BodyAdjacentTracePair}
  v\in L^2(0,T;H^{s+1}),
  \qquad
  v_t\in L^2(0,T;H^{s-1}).
\end{equation}
Cauchy--Schwarz gives
\begin{equation}\label{eq:BodyAdjacentTraceEstimate}
  \left|
  \norm{v(b)}_{H^s}^2-
  \norm{v(a)}_{H^s}^2
  \right|
  \leq
  2\norm{v_t}_{L^2(a,b;H^{s-1})}
  \norm{v}_{L^2(a,b;H^{s+1})}.
\end{equation}
Time mollification and Fourier truncation extend the identity to the pair in \eqref{eq:BodyAdjacentTracePair}.
Weak \(H^s\) continuity together with continuity of the norm gives a unique representative
\begin{equation}\label{eq:BodyAdjacentTraceConclusion}
  v\in C([0,T];H^s).
\end{equation}

The first continuation-grade example uses \(v=u\) and the middle height \(s=2\):
\begin{equation}\label{eq:FirstConcreteAdjacentPair}
  u\in L^2(0,T;H^3),
  \qquad
  u_t\in L^2(0,T;H^1)
  \quad\Longrightarrow\quad
  u\in C([0,T];H^2).
\end{equation}
The two spaces are offset by two spatial derivatives because they meet at the middle space \(H^2\).
The upper row controls the field one derivative above the trace height.
The lower row controls its time derivative one derivative below.

Apply the same pair to \(v=V_n\), whose time derivative is \(V_{n+1}\).
Fix finite temporal depth \(N\) and middle Sobolev order \(M\).
The velocity size of the corresponding finite block is
\begin{equation}\label{eq:BodyRectangleNorm}
\begin{aligned}
  \mathfrak R_{N,M}(u;T)
  :={}&
  \sum_{n=0}^{N}
  \norm{V_n}_{L^2(0,T;H^{M+1})}^2
  \\
  &+
  \sum_{n=0}^{N}
  \norm{V_{n+1}}_{L^2(0,T;H^{M-1})}^2.
\end{aligned}
\end{equation}
When \(M=0\), the lower edge is the inhomogeneous space \(H^{-1}\).
No one constant across all values of \(N\) and \(M\) is part of this definition.

At \((N,M)=(1,2)\), the rectangle retains
\[
  u,u_t\in L^2(0,T;H^3),
  \qquad
  u_t,u_{tt}\in L^2(0,T;H^1).
\]
The repeated \(u_t\) is not duplicated bookkeeping.
It is the lower member of the pair that traces \(u\) and the upper member of the pair that traces \(u_t\).
The temporal pairs overlap on one actual derivative of the same solution.

The pressure-complete finite rectangle \(\cR_{N,M}[u_0;T]\) retains these velocity rows, the pressure rows \(Q_0,\ldots,Q_N\), and the recurrences that identify every entry with the mixed jet generated before the cutoffs were chosen.
For a strict cutoff \(T<T_*\), classicality on the compact interval \([0,T]\) gives
\begin{equation}\label{eq:StrictSubintervalRectangleBound}
  \mathfrak R_{N,M}(u;T)<\infty
  \qquad\text{for every finite }N,M.
\end{equation}
That statement does not reach the candidate endpoint.

The distinction is a quantifier, not a technicality.
The scalar function \(f(t)=(T_*-t)^{-1/2}\) belongs to \(L^2(0,T)\) for every \(T<T_*\), but not to \(L^2(0,T_*)\).
This is an integrability countertest, not a fluid scenario.
It shows why finite control after the cutoff has been chosen does not produce one bound that survives every cutoff.

The exact unpaid statement is now visible.
After \(N\) and \(M\) have been fixed, the datum and viscosity must control the corresponding pressure-complete rectangle with one constant independent of \(T<T_*\).
That whole-terminal estimate is the central theorem of the next section.
